\section{Introduction}
\IEEEPARstart{F}{ederated} learning (FL) trains a model across many clients---phones, hospitals, banks---that never
send their raw data to a server. Each client improves the current model on its own records and returns only the
resulting update, which the server averages into the next model~\cite{mcmahan2017communication,li2020federated}.
This makes collaboration possible where pooling data is forbidden by law or by commercial interest. Its central
difficulty is heterogeneity: client datasets are not identically distributed, so the updates being averaged pull
in different directions~\cite{kairouz2021advances}.

Keeping data local is not the same as keeping it private. A shared gradient can be inverted to reconstruct the
training examples behind it~\cite{zhu2019deep}, and a released model can reveal whether a particular record was
used to train it~\cite{shokri2017membership,nasr2019comprehensive}. The accepted defence is differential privacy
(DP)~\cite{dwork2006calibrating}, usually applied through differentially private stochastic gradient descent
(DP-SGD)~\cite{abadi2016deep}: each example's gradient is clipped to a fixed norm, Gaussian noise is added, and an
accountant tracks the cumulative privacy loss as a budget~$\eps$.

DP has a side effect that matters for fairness. Bagdasaryan \emph{et al.}~\cite{bagdasaryan2019differential} showed
that DP-SGD does not lower accuracy evenly: groups the model already serves poorly lose the most, because their
examples produce large gradients that clipping truncates hardest. Esipova \emph{et al.}~\cite{esipova2023disparate}
traced the effect further, to the misalignment between clipped and true gradient directions. In FL the ``groups''
are often the clients themselves, so a privacy guarantee can quietly widen the gap between well-served and
poorly-served participants~\cite{chen2023privacy,rafi2024fairness}.

Personalisation is a natural counterweight. Instead of forcing every client onto one global compromise,
personalised FL gives each client a model adapted to its own data~\cite{tan2023towards}, and Bietti
\emph{et al.}~\cite{bietti2022personalization} proved that personalised models tolerate privacy noise better than a
single global model. Personalisation, however, is not one technique but several families: local fine-tuning of the
global model, meta-learning of an easily adapted starting point~\cite{fallah2020personalized}, clustering of
similar clients~\cite{ghosh2022efficient}, and decoupling of shared and private
parameters~\cite{arivazhagan2019federated,collins2021exploiting}. Each family was introduced in its own study, with
its own datasets, baselines and hyperparameters, and in most cases without any privacy mechanism. The closest recent
work, DP-Ditto~\cite{zhao2025convergence}, extends one personalisation method under DP and compares other methods
against it; it establishes the value of that method but is not a neutral comparison of families. A practitioner who
must deploy at a given~$\eps$ therefore has no principled way to choose among them.

This paper closes that gap with a controlled comparison. We hold the data partition, model, clipping norm, sampling
rate, number of rounds and privacy accounting fixed, and change only the personalisation mechanism. Privacy is
protected at the level of individual records inside each client, and every mechanism is charged for every query it
makes on private data, so all of them spend the same $(\eps,\delta)$. The contributions are:
\begin{itemize}
  \item \textbf{An equal-budget protocol.} Every client answers the same number of noisy gradient queries; the loss
  values that clustering and q-FFL need are themselves privatised and charged to the same accountant; and the noise
  multiplier of each run is calibrated so that every mechanism ends at the same $(\eps,\delta)$.
  \item \textbf{A privacy sweep} from $\eps=0.3$ to $\eps=8$ (and the non-private case) that measures how DP
  redistributes accuracy across 20 heterogeneous clients and tests whether there is a budget below which
  personalisation stops paying off.
  \item \textbf{A head-to-head comparison} of fine-tuning, Per-FedAvg, IFCA-style clustering and FedPer against
  DP-FedAvg and local-only training, scored by mean and worst-10\% client accuracy, Jain's fairness index,
  personalisation gain, group-fairness gaps, measured membership-inference leakage, and communication and compute
  cost---summarised as a recommended operating configuration for each budget.
  \item \textbf{A fairness-objective test}: whether q-FFL~\cite{li2020fair} raises worst-client accuracy at a fixed
  $\eps$, and whether the gain is paid for in mean accuracy or in privacy budget.
\end{itemize}

Section~\ref{sec:related} reviews related work and states the gap. Section~\ref{sec:method} formalises the problem
and the mechanisms. Section~\ref{sec:setup} describes the experimental setup, Section~\ref{sec:results} presents the
results, Section~\ref{sec:discussion} discusses them together with their limitations, and
Section~\ref{sec:conclusion} concludes.
